% GENERATED FILE. Edit canonical lessons, then run npm run generate:books.
% canonical-source-hash: fnv1a64:cf943fdbfad2fe68
% canonical-lessons: SA-C374-avahitah, SA-C374-pramadah, SA-C374-pariharati, SA-C374-anarthah, SA-C374-viramati

\chapter{Taking Care}
\label{ch:sa-taking-care}
\hlchaptermodality{374}

\begin{chapteropening}
\textbf{By the end of this chapter:} \emph{I can warn someone what will happen if they are careless.}

Complete the last lesson of chapter 374: I can warn someone what will happen if they are careless.
\end{chapteropening}

\section[\texorpdfstring{\sk{अवहितः} (avahitaḥ)}{अवहितः (avahitaḥ)}]{\sk{अवहितः} (avahitaḥ) — attentive; careful}

\label{lesson:SA-C374-avahitah}



\begin{quote}
Before the new one: say the Sanskrit for an exception, then the Sanskrit for compulsory.
\end{quote}

\subsection*{You'll want to know: \sk{अवहितः}}
\textbf{\sk{अवहितः}} — \emph{avahitaḥ} — \textquotedblleft{}attentive; careful\textquotedblright{}.

\begin{quote}
\sk{यदि} \sk{भवान्} \sk{अवहितः}, \sk{तर्हि} \sk{दोषः} \sk{न} \sk{भविष्यति।}
\end{quote}

\begin{quote}
\emph{yadi bhavān avahitaḥ, tarhi doṣaḥ na bhaviṣyati.}
\end{quote}

\begin{quote}
If you are attentive, there will be no mistake.
\end{quote}

\begin{cousinweb}
\textbf{\sk{अवहितः}} (\emph{avahitaḥ}) is built on \emph{ava}, \textquotedblleft{}down\textquotedblright{}, and \emph{hita}, \textquotedblleft{}put\textquotedblright{}, the participle of √dhā; MW gives \emph{ava-hita} \textquotedblleft{}attentive\textquotedblright{} (MW).
\end{cousinweb}

\begin{grammarlens}[title={What you've built}]
One more word for saying what happens if.
\end{grammarlens}

\subsection*{Your turn}
\begin{itemize}
\raggedright
  \item \emph{Say it:} \emph{avahitaḥ}
  \item \emph{Say it:} \emph{avahitaḥ}, once more
  \item \emph{Say it:} \emph{yadi bhavān avahitaḥ, tarhi doṣaḥ na bhaviṣyati.} — If you are attentive, there will be no mistake.
  \item \emph{Recall:} say the Sanskrit for an exception, then the Sanskrit for compulsory, then say \emph{avahitaḥ} again
\end{itemize}

\begin{tcolorbox}[breakable,colback=teal!4,colframe=teal!35!black,title={Before you move on}]
What does \sk{अवहितः} mean? (\textquotedblleft{}Attentive; careful\textquotedblright{}.) Say it once more.
\end{tcolorbox}
\section[\texorpdfstring{\sk{प्रमादः} (pramādaḥ)}{प्रमादः (pramādaḥ)}]{\sk{प्रमादः} (pramādaḥ) — carelessness}

\label{lesson:SA-C374-pramadah}



\begin{quote}
Before the new one: say the Sanskrit for compulsory, then the Sanskrit for attentive.
\end{quote}

\subsection*{You'll want to know: \sk{प्रमादः}}
\textbf{\sk{प्रमादः}} — \emph{pramādaḥ} — \textquotedblleft{}carelessness\textquotedblright{}.

\begin{quote}
\sk{यदि} \sk{प्रमादः} \sk{भवति}, \sk{तर्हि} \sk{हानिः} \sk{भविष्यति।}
\end{quote}

\begin{quote}
\emph{yadi pramādaḥ bhavati, tarhi hāniḥ bhaviṣyati.}
\end{quote}

\begin{quote}
If there is carelessness, there will be a loss.
\end{quote}

A masculine noun.

\begin{cousinweb}
\textbf{\sk{प्रमादः}} (\emph{pramādaḥ}) is built on \emph{pra} and the root √mad, \textquotedblleft{}to be drunk\textquotedblright{}; MW gives \emph{pra-māda} \textquotedblleft{}negligence, carelessness\textquotedblright{} (MW).
\end{cousinweb}

\begin{grammarlens}[title={What you've built}]
One more word for saying what happens if.
\end{grammarlens}

\subsection*{Your turn}
\begin{itemize}
\raggedright
  \item \emph{Say it:} \emph{pramādaḥ}
  \item \emph{Say it:} \emph{pramādaḥ}, once more
  \item \emph{Say it:} \emph{yadi pramādaḥ bhavati, tarhi hāniḥ bhaviṣyati.} — If there is carelessness, there will be a loss.
  \item \emph{Recall:} say the Sanskrit for compulsory, then the Sanskrit for attentive, then say \emph{pramādaḥ} again
\end{itemize}

\begin{tcolorbox}[breakable,colback=teal!4,colframe=teal!35!black,title={Before you move on}]
What does \sk{प्रमादः} mean? (\textquotedblleft{}Carelessness\textquotedblright{}.) Say it once more.
\end{tcolorbox}
\section[\texorpdfstring{\sk{परिहरति} (pariharati)}{परिहरति (pariharati)}]{\sk{परिहरति} (pariharati) — avoids}

\label{lesson:SA-C374-pariharati}



\begin{quote}
Before the new one: say the Sanskrit for attentive, then the Sanskrit for carelessness.
\end{quote}

\subsection*{You'll want to know: \sk{परिहरति}}
\textbf{\sk{परिहरति}} — \emph{pariharati} — \textquotedblleft{}avoids\textquotedblright{}.

\begin{quote}
\sk{यदि} \sk{भवान्} \sk{क्रोधं} \sk{परिहरति}, \sk{तर्हि} \sk{सुखी} \sk{भविष्यति।}
\end{quote}

\begin{quote}
\emph{yadi bhavān krodhaṃ pariharati, tarhi sukhī bhaviṣyati.}
\end{quote}

\begin{quote}
If you avoid anger, you will be happy.
\end{quote}

Past (he …): \sk{सः} \sk{विवादं} \sk{परिहृतवान्}, \emph{saḥ vivādaṃ parihṛtavān}.

\begin{cousinweb}
\textbf{\sk{परिहरति}} (\emph{pariharati}) is built on \emph{pari}, \textquotedblleft{}around\textquotedblright{}, and the root √hṛ, \textquotedblleft{}to take\textquotedblright{}; MW gives \emph{pari-hṛta} \textquotedblleft{}shunned, avoided\textquotedblright{} (MW).
\end{cousinweb}

\begin{grammarlens}[title={What you've built}]
One more word for saying what happens if.
\end{grammarlens}

\subsection*{Your turn}
\begin{itemize}
\raggedright
  \item \emph{Say it:} \emph{pariharati}
  \item \emph{Say it:} \emph{pariharati}, once more
  \item \emph{Say it:} \emph{yadi bhavān krodhaṃ pariharati, tarhi sukhī bhaviṣyati.} — If you avoid anger, you will be happy.
  \item \emph{Recall:} say the Sanskrit for attentive, then the Sanskrit for carelessness, then say \emph{pariharati} again
\end{itemize}

\begin{tcolorbox}[breakable,colback=teal!4,colframe=teal!35!black,title={Before you move on}]
What does \sk{परिहरति} mean? (\textquotedblleft{}Avoids\textquotedblright{}.) Say it once more.
\end{tcolorbox}
\section[\texorpdfstring{\sk{अनर्थः} (anarthaḥ)}{अनर्थः (anarthaḥ)}]{\sk{अनर्थः} (anarthaḥ) — a mishap; a bad turn}

\label{lesson:SA-C374-anarthah}



\begin{quote}
Before the new one: say the Sanskrit for carelessness, then the Sanskrit for avoids.
\end{quote}

\subsection*{You'll want to know: \sk{अनर्थः}}
\textbf{\sk{अनर्थः}} — \emph{anarthaḥ} — \textquotedblleft{}a mishap; a bad turn\textquotedblright{}.

\begin{quote}
\sk{यदि} \sk{चालकः} \sk{श्रान्तः}, \sk{तर्हि} \sk{अनर्थः} \sk{भविष्यति।}
\end{quote}

\begin{quote}
\emph{yadi cālakaḥ śrāntaḥ, tarhi anarthaḥ bhaviṣyati.}
\end{quote}

\begin{quote}
If the driver is tired, there will be a mishap.
\end{quote}

A masculine noun.

\begin{cousinweb}
\textbf{\sk{अनर्थः}} (\emph{anarthaḥ}) is built from \emph{an-}, \textquotedblleft{}not\textquotedblright{}, and \emph{artha}, \textquotedblleft{}aim, use\textquotedblright{}; MW gives \emph{an-artha} \textquotedblleft{}a worthless object\textquotedblright{} and \textquotedblleft{}a disappointing occurrence, a reverse, an evil\textquotedblright{} (MW).
\end{cousinweb}

\begin{grammarlens}[title={What you've built}]
One more word for saying what happens if.
\end{grammarlens}

\subsection*{Your turn}
\begin{itemize}
\raggedright
  \item \emph{Say it:} \emph{anarthaḥ}
  \item \emph{Say it:} \emph{anarthaḥ}, once more
  \item \emph{Say it:} \emph{yadi cālakaḥ śrāntaḥ, tarhi anarthaḥ bhaviṣyati.} — If the driver is tired, there will be a mishap.
  \item \emph{Recall:} say the Sanskrit for carelessness, then the Sanskrit for avoids, then say \emph{anarthaḥ} again
\end{itemize}

\begin{tcolorbox}[breakable,colback=teal!4,colframe=teal!35!black,title={Before you move on}]
What does \sk{अनर्थः} mean? (\textquotedblleft{}A mishap; a bad turn\textquotedblright{}.) Say it once more.
\end{tcolorbox}
\section[\texorpdfstring{\sk{विरमति} (viramati)}{विरमति (viramati)}]{\sk{विरमति} (viramati) — ceases; stops}

\label{lesson:SA-C374-viramati}



\begin{quote}
Before the new one: say the Sanskrit for avoids, then the Sanskrit for a mishap.
\end{quote}

\subsection*{You'll want to know: \sk{विरमति}}
\textbf{\sk{विरमति}} — \emph{viramati} — \textquotedblleft{}ceases; stops\textquotedblright{}.

\begin{quote}
\sk{यदि} \sk{वर्षा} \sk{विरमति}, \sk{तर्हि} \sk{अहं} \sk{बहिः} \sk{गमिष्यामि।}
\end{quote}

\begin{quote}
\emph{yadi varṣā viramati, tarhi ahaṃ bahiḥ gamiṣyāmi.}
\end{quote}

\begin{quote}
If the rain stops, I will go out.
\end{quote}

Past (it …): \sk{वर्षा} \sk{विरतवती}, \emph{varṣā viratavatī}.

\begin{cousinweb}
\textbf{\sk{विरमति}} (\emph{viramati}) is built on \emph{vi} and the root √ram, \textquotedblleft{}to rest\textquotedblright{}; MW gives \emph{vi-√ram} \textquotedblleft{}to stop, to pause, to cease, to come to an end\textquotedblright{} (MW).
\end{cousinweb}

\begin{grammarlens}[title={What you've built}]
One more word for saying what happens if.
\end{grammarlens}

\subsection*{Your turn}
\begin{itemize}
\raggedright
  \item \emph{Say it:} \emph{viramati}
  \item \emph{Say it:} \emph{viramati}, once more
  \item \emph{Say it:} \emph{yadi varṣā viramati, tarhi ahaṃ bahiḥ gamiṣyāmi.} — If the rain stops, I will go out.
  \item \emph{Recall:} say the Sanskrit for avoids, then the Sanskrit for a mishap, then say \emph{viramati} again
\end{itemize}

\begin{tcolorbox}[breakable,colback=teal!4,colframe=teal!35!black,title={Before you move on}]
What does \sk{विरमति} mean? (\textquotedblleft{}Ceases; stops\textquotedblright{}.) Say it once more.
\end{tcolorbox}
